\section{Raw sources：互联网不是一个可直接训练的数据集}

“语言模型训练在整个互联网”是非常粗糙的说法。更准确地说，很多原始数据来自 public web 的 crawls，但 live servers 不能直接训练。需要 crawler 发现 URL、下载页面、保存 raw HTTP response，再把 HTML 等格式转换为文本。

\begin{knowledgebox}{课堂提示：不要把“整个互联网”当作可训练集合}
老师先用这个常见说法制造反差，再逐项拆出技术与法律限制：动态页面和登录墙可能抓不到，robots.txt、rate limit 与 Terms of Service 会限制访问，版权和隐私又会限制用途。因而“在网上可见”只描述了潜在来源，并没有说明数据是否可达、可解析、可授权或适合训练。
\end{knowledgebox}

\begin{knowledgebox}{术语消化：从 live service 到 processed data}
\begin{tabular}{L{0.22\textwidth}L{0.34\textwidth}L{0.34\textwidth}}
\toprule
层次 & 例子 & 问题 \\
\midrule
Live service & 网站、GitHub、arXiv、Wikipedia。 & 动态页面、认证、rate limit、ToS、robots.txt。 \\
Raw data & WARC、git repo、PDF、LaTeX source、XML dump。 & 格式复杂，噪声大，重复多。 \\
Processed data & WET/text、filtered docs、deduped corpus。 & 转换损失、过滤 bias、PII、版权和质量控制。 \\
\bottomrule
\end{tabular}
\end{knowledgebox}

\begin{figure}[H]
\centering
\includegraphics[width=0.82\textwidth]{decline-consent.png}
\caption{Decline of consent：C4、RefinedWeb、Dolma 等常见数据 URL 的 robots.txt/ToS 限制随时间增加。}
\figsource{CS336 Lecture 13 官方源引用图。}
\end{figure}

\begin{knowledgebox}{读图：限制增加意味着什么}
越来越多网站通过 robots.txt、ToS 或其他方式表达不希望被爬取。即使技术上可访问，也不代表合规、礼貌或可用于训练。未来数据获取的成本会更高，许可数据和自有数据的重要性会上升。
\end{knowledgebox}

\begin{figure}[H]
\centering
\includegraphics[width=0.55\textwidth]{anthropic-crawling.png}
\caption{不良 crawler 行为会造成服务压力和公众反弹。}
\figsource{CS336 Lecture 13 官方源引用图。}
\end{figure}

\begin{warningbox}{Shadow libraries 的边界}
LibGen、Z-Library、Anna's Archive、Sci-Hub 等 technically 是 web 的一部分，但它们绕过版权和付费墙。课程强调从法律角度这是 piracy/copyright infringement。训练数据工程不能把“能下载”当成“能使用”。
\end{warningbox}

\begin{importantbox}{本章小结：raw source 不是 raw permission}
网页能访问、能下载、能解析，并不意味着能训练。真实数据管线至少要同时处理四种约束：技术可达性、站点意愿、法律授权、社会接受度。这四个维度任一出问题，都可能让数据集不可持续。
\end{importantbox}

